%================================
% Vocabulary & Language Focus — XB2 (Unit 1)
% 编译方式: xelatex main.tex (需运行两遍以生成目录)
%================================
\documentclass[../main.tex]{subfiles}

\begin{document}

\section{Unit 1}

\vw{critical}{/ˈkrɪtɪkl/}{adj.}

\sense{adj.}{批评的，批判性的}
\col{be critical of sb/sth（对……挑剔、批评）}
\ex{The report was highly critical of the government's handling of the crisis.}
\col{critical thinking 批判性思维}
\ex{Critical thinking means questioning what you read rather than accepting it blindly.}

\sense{adj.}{关键的}
\col{be critical to/for sth}
\ex{The choice of location is critical to the success of the project.}
\col{a critical moment/point}
\ex{We are at a critical moment in the negotiations.}

\sense{adj.}{危急的}
\col{be in a critical condition}
\ex{Two of the injured are still in a critical condition.}

\begin{origin}
来自希腊语 \textit{kritikos}（有判断力的），源于动词 \textit{krinein}（判断、区分、裁决）；经由拉丁语 \textit{criticus}、法语 \textit{critique} 进入英语。同源词 \textbf{crisis}（危机）本义是"必须做出裁决的紧要关头"，\textbf{criterion}（标准）是"用来判断的尺子"。批评、关键时刻与危急状况，共享同一个"判断"意象。
\end{origin}

\begin{lex}{critical / crucial / vital}
三者都可表"至关重要"，常可互换，但侧重不同：\textbf{critical} 强调"成败所系"（\textit{a critical decision} 决定成败的一步）；\textbf{crucial} 强调"决定性"（\textit{a crucial experiment} 决定性实验）；\textbf{vital} 本义"与生命攸关"（\textit{vital organs} 要害器官），口语中语气最重。
\end{lex}

\vw{trap}{/træp/}{vt. \& n.}

\sense{vt.}{使落入险境}
\col{be trapped in/under/behind sth}
\ex{The miners were trapped underground for three days.}
\col{be trapped by the fire/flood}
\ex{Rescuers saved a family trapped by the fire.}

\sense{vt.}{卡住，绊住}
\col{get trapped in/under sth}
\ex{She got her finger trapped in the door.}
\ex{He trapped his hand in the car door.}

\sense{n.}{陷阱，罗网；圈套}
\col{set/lay a trap (for sb)}
\ex{The police set a trap to catch the burglar.}
\col{fall/walk into a trap}
\ex{Don't fall into the trap of believing everything you read online.}

\sense{n.}{困境，牢笼}
\col{a trap of sth}
\ex{Millions of families are caught in the trap of poverty.}
\col{feel trapped（感到身陷牢笼）}
\ex{She felt trapped in a loveless marriage.}

\begin{origin}
源自古英语 \textit{træppe}（罗网），与原始日耳曼语 *\textit{trap-}（踩踏）有关——陷阱本指"踩上去就会翻倒的机关"；同源词如中世纪低地德语 \textit{trappe}（楼梯，即"踏脚的东西"）。原始印欧语词根 *\textit{dreb-} 意为"行走、踩踏"。
\end{origin}

\vw{release}{/rɪˈliːs/}{vt. \& n.}

\sense{vt.}{发布}
\col{release a film/album/product}
\ex{The new album will be released next month.}
\col{a press release 新闻稿}
\ex{The company issued a press release about the merger.}

\sense{vt.}{释放}
\col{release sb from sth}
\ex{The hostages were released unharmed.}
\ex{The nurse was released from her duties for the day.}

\sense{vt.}{松开}
\col{release the brake/clutch/hand}
\ex{Release the handbrake before you drive off.}

\sense{vt.}{发泄}
\col{release tension/anger/energy}
\ex{Jogging is a good way to release stress.}

\sense{n.}{释放；发行；排放，泄漏}
\col{the release of sb/sth}
\ex{the release of the prisoners}
\ex{the release of toxic gases into the river}

\begin{origin}
来自古法语 \textit{relessier}（现代法语 \textit{relâcher}），最终源自拉丁语 \textit{relaxare}（放松、松开），即 \textit{re-}（回）+ \textit{laxare}（松开），词根 \textit{laxus}（松弛的）。因此 \textbf{release} 与 \textbf{relax} 是同一拉丁词衍生出的"双重词"（doublet）——"释放"和"放松"本是同一个动作：松开手。
\end{origin}

\vw{extend}{/ɪkˈstend/}{vi. \& vt.}

\sense{vt.}{延长}
\col{extend a deadline/visa/period}
\ex{The government extended the deadline by two weeks.}
\ex{She applied to have her visa extended.}

\sense{vt.}{扩大}
\col{extend a building/road}
\ex{They plan to extend the kitchen to make more room.}

\sense{vt.}{伸展}
\col{extend one's hand/arm}
\ex{She extended her hand in greeting.}

\sense{vt.}{提供，表示（正式）}
\col{extend sth to sb}
\ex{We would like to extend a warm welcome to our new members.}
\ex{The bank refused to extend further credit to the firm.}

\sense{vi.}{延伸（距离）}
\col{extend for/from... to...}
\ex{The forest extends for miles along the coast.}
\col{extend over（持续、覆盖）}
\ex{The hot weather extended over several weeks.}

\begin{origin}
来自拉丁语 \textit{extendere}（向外伸展）：\textit{ex-}（向外）+ \textit{tendere}（拉伸）。同源词：\textbf{tension}（张力）、\textbf{tendon}（肌腱）、\textbf{extensive}（广阔的）、\textbf{extent}（程度、范围）。"tendere 拉伸"是英语中极其多产的一支词根。
\end{origin}

\begin{lex}{extend / expand}
\textbf{extend} 侧重长度、时间、范围上的"延伸"（\textit{extend the deadline} 延长期限）；\textbf{expand} 侧重体积、规模上的"扩张"（\textit{metal expands when heated} 金属受热膨胀；\textit{expand the business} 扩张业务）。前者像拉橡皮筋，后者像吹气球。
\end{lex}

\vw{construction}{/kənˈstrʌkʃn/}{n.}

\sense{n.}{建筑，建造}
\col{under construction 在建造中}
\ex{The new bridge is still under construction.}
\col{the construction of sth}
\ex{The construction of the dam took nearly ten years.}
\col{a construction site 建筑工地}
\ex{Workers at the construction site must wear helmets.}

\sense{n.}{建造物}
\ex{The tower is a modern construction of glass and steel.}

\begin{origin}
来自拉丁语 \textit{constructio}，源于 \textit{construere}（堆叠建造）：\textit{con-}（一起）+ \textit{struere}（堆、筑）。同源词：\textbf{structure}（结构）、\textbf{instruct}（in- 里 + struere 搭建 $\to$ 把知识搭进头脑 $\to$ 教导）、\textbf{destroy}（de- 拆掉 + struere $\to$ 把堆起来的东西拆散）。
\end{origin}

\vw{minor}{/ˈmaɪnə(r)/}{adj.}

\sense{adj.}{轻微的}
\col{a minor injury/operation/repair}
\ex{He suffered only minor injuries in the accident.}
\ex{The car needs a few minor repairs.}

\sense{adj.}{次要的}
\col{a minor role/issue/character}
\ex{She played only a minor role in the negotiations.}

\sense{n.}{未成年人（拓展）}
\ex{Minors are not allowed to buy alcohol in most countries.}
\col{minor in sth（大学）辅修}
\ex{She minored in psychology at college.}

\begin{origin}
直接借自拉丁语 \textit{minor}（更小的）——它是拉丁语形容词 \textit{parvus}（小）的比较级，英语干脆把整个比较级借了过来。同源词：\textbf{minus}（减去）、\textbf{minority}（少数）、\textbf{minuscule}（极小的）。反义词 \textbf{major} 同样是拉丁语比较级 \textit{maior}（更大的）。拉丁语比较级词尾 -or 原样进入了英语。
\end{origin}

\vw{scream}{/skriːm/}{vi. \& vt. \& n.}

\sense{vi.}{尖叫}
\col{scream with/in sth}
\ex{The fans screamed with excitement when the star appeared.}
\col{scream at sb}
\ex{She screamed at the children to stop fighting.}

\sense{vi.}{高声喊}
\col{scream for help}
\ex{We could hear someone screaming for help in the distance.}

\sense{vi.}{发出大而尖的声音}
\ex{The brakes screamed as the car stopped suddenly.}

\sense{n.}{尖叫，尖锐刺耳的声音}
\col{a scream of delight/terror}
\ex{She let out a scream of delight.}
\ex{We heard a blood-curdling scream from the next room.}

\begin{origin}
中古英语 \textit{scremen}，可能借自古诺斯语 \textit{skræma}（吓唬、尖叫）。日耳曼语族里有一族以 skr-/shr- 开头的拟声词，模拟尖锐刺耳的声响：\textbf{shriek}（尖叫）、\textbf{screech}（尖啸）、\textbf{creak}（吱嘎作响）。scream 本就是一出生就在模仿自己的发音。
\end{origin}

\vw{bark}{/bɑːk/}{vi. \& vt. \& n.}

\sense{vi.}{（狗）吠叫}
\col{bark at sb/sth}
\ex{The dog barked furiously at the postman.}

\sense{vt.}{厉声发令}
\col{bark out orders}
\ex{He barked out orders to his assistants.}
\col{bark at sb}
\ex{"Don't interrupt!" she barked at me.}

\sense{n.}{（狗）吠声；短促响亮的人声}
\col{a bark of laughter}
\ex{He gave a short bark of laughter.}
\col{His bark is worse than his bite. 刀子嘴，豆腐心}
\ex{Don't be afraid of the boss — his bark is worse than his bite.}

\sense{n.}{树皮}
\col{the bark of a tree}
\ex{The bark of this tree peels off in thin strips.}
\col{bark up the wrong tree 找错对象，白费力气}
\ex{If you think I stole the money, you're barking up the wrong tree.}

\begin{origin}
英语里罕见的"同形异源"词：表示"吠叫"的 \textit{bark} 源自古英语 \textit{beorcan}，属拟声词（模拟狗叫）；表示"树皮"的 \textit{bark} 则借自北欧语（如古诺斯语 \textit{bǫrkr}）。两个词拼写相同、词源毫不相干，属于同形异义词（homograph）。
\end{origin}

\vw{choke}{/tʃəʊk/}{vt. \& vi.}

\sense{vt. \& vi.}{（使）窒息}
\col{choke on sth}
\ex{He almost choked on a fish bone.}
\col{be choked by smoke}
\ex{The rescuers were nearly choked by the smoke.}
\col{choke to death}
\ex{The cat choked to death on a bone.}

\sense{vt. \& vi.}{（使）哽咽，（尤指感情激动而）说不出话来}
\col{choke with emotion/anger}
\ex{Her voice choked with emotion as she thanked the rescuers.}
\col{choke back tears 强忍住眼泪}
\ex{She choked back tears as she read the letter.}

\col{be choked with sth 被……堵塞}
\ex{The streets were choked with traffic.}

\begin{origin}
中古英语 \textit{choken}，源自古英语 \textit{ācēocian}（使窒息），与 \textit{cēoce}（颚、面颊）同根——本义是"卡在颚部、喉咙里"。同源词正是 \textbf{cheek}（面颊）。原来"窒息"和"脸颊"共享同一个身体意象。
\end{origin}

\vw{carpet}{/ˈkɑːpɪt/}{n.}

\sense{n.}{地毯}
\col{lay a carpet}
\ex{We laid a new carpet in the living room.}
\col{roll out the red carpet 隆重欢迎}
\ex{The city rolled out the red carpet for the visiting delegation.}

\sense{n.}{覆盖地面的一层厚东西}
\col{a carpet of sth}
\ex{The ground was covered with a carpet of fallen leaves.}
\ex{A thick carpet of snow lay over the fields.}

\col{on the carpet 受训斥（口语）}
\ex{The manager had him on the carpet for being late again.}

\begin{origin}
来自古法语 \textit{carpite}（厚织的装饰布），源于拉丁语 \textit{carpita}（vestis），即"梳理过的厚羊毛布"——词根 \textit{carpere}（采摘、梳理羊毛）。同源词：\textbf{excerpt}（ex- 摘出 + carpere 摘 $\to$ 摘录）。毛毡制品须先"梳理"纤维，地毯由此得名。
\end{origin}

\vw{automatic}{/ˌɔːtəˈmætɪk/}{adj.}

\sense{adj.}{自动的}
\col{automatic doors/transmission}
\ex{The building has automatic doors at every entrance.}
\col{an automatic rifle/weapon 自动步枪}

\sense{adj.}{无意识的，不假思索的}
\col{an automatic response/reaction}
\ex{Thanking people has become an automatic response for her.}
\ex{Breathing is automatic — we never think about it.}

\begin{origin}
来自希腊语 \textit{automatos}（自行移动的）：\textit{autos}（自己）+ \textit{matos}（思考、行动的，源自 \textit{menos} 意愿、动力）。同源词：\textbf{automobile}（automobile：自己会走）、\textbf{automation}（自动化）。"自己会动"——两千多年前希腊人已造出这个词，只是当时没有机器可形容。
\end{origin}

\vw{investigate}{/ɪnˈvestɪɡeɪt/}{vt. \& vi.}

\sense{vt. \& vi.}{侦查，调查}
\col{investigate a crime/complaint/claim}
\ex{Police are investigating the cause of the fire.}
\col{investigate whether/why/how...}
\ex{Scientists are investigating why the fish are dying.}

\sense{vt.}{研究}
\col{investigate the effects/causes of sth}
\ex{Researchers investigated the effects of sleep on memory.}

\col{investigate + 宾语（不接 that 从句）}
\ex{The committee is investigating the allegations of corruption.}

\begin{origin}
来自拉丁语 \textit{investigare}（循迹追查）：\textit{in-}（进入）+ \textit{vestigium}（脚印、足迹）。同源词：\textbf{vestige}（遗迹、残余）、\textbf{vestigial}（退化的）。调查本义是"顺着脚印一路追踪"——侦探和科学家做的是同一件事。
\end{origin}

\vp{come about}

\sense{v.}{发生}
\col{How did it come about that ...?（正式）}
\ex{How did it come about that the two rivals became partners?}
\ex{The discovery came about quite by chance.}
\ex{It came about that the two brothers met again after twenty years.}

\col{比 happen/occur 更书面，常用于 how 问句与巧合叙述}

\begin{origin}
字面义为"转着到来、绕着来到"：come（来）+ about（周围）。16 世纪起由"绕过来"引申出"发生"。与之同构的 \textbf{happen} 则源自北欧语 \textit{hap}（运气）——"发生的事"在英语词源里要么是"绕过来的"，要么是"碰运气的"。
\end{origin}

\vw{priority}{/praɪˈɒrəti/}{n.}

\sense{n.}{优先事项，首要事情}
\col{a top/first priority}
\ex{Education remains a top priority for the government.}
\col{get one's priorities right/straight 分清轻重缓急}
\ex{You need to get your priorities right before it's too late.}

\sense{n.}{优先权，重点}
\col{give priority to sth}
\ex{Give priority to the most urgent cases first.}
\col{take/have priority over sth}
\ex{Safety must take priority over profit.}

\begin{origin}
来自拉丁语 \textit{prior}（更靠前的）——\textit{prae}（在前）的比较级——再加名词后缀 -ity。字面义即"排在其他事情前面的事"。同源词：\textbf{prior}（先前的；prior knowledge 已有知识）。
\end{origin}

\vw{contradict}{/ˌkɒntrəˈdɪkt/}{vt.}

\sense{vt.}{相矛盾，相反}
\col{contradict sth}
\ex{The two statements contradict each other.}
\ex{These results contradict earlier findings.}

\sense{vt.}{反驳，驳斥}
\col{contradict sb}
\ex{Don't contradict your father in front of the guests.}
\col{contradict oneself 自相矛盾}
\ex{He contradicted himself twice during the speech.}

\begin{origin}
来自拉丁语 \textit{contradicere}（对着说、反驳）：\textit{contra}（相反）+ \textit{dicere}（说）。同源词：\textbf{dictate}（口述、命令）、\textbf{predict}（pre- 预先 + dict 说 $\to$ 预言）、\textbf{dictionary}（diction 措辞 + -ary $\to$ 汇集言词的书）。"话对着话"就是矛盾与反驳。
\end{origin}

\vw{factual}{/ˈfæktʃuəl/}{adj.}

\sense{adj.}{事实的，真实的}
\col{factual information/errors}
\ex{The article is full of factual errors.}
\col{a factual account/report}
\ex{The police gave a factual account of what happened.}
\col{factual accuracy}
\ex{We must check the factual accuracy of the story.}

\begin{origin}
来自拉丁语 \textit{factum}（做出来的事）——动词 \textit{facere}（做）的过去分词名词化；加形容词后缀 -al。同源词：\textbf{fact}（事实）、\textbf{factory}（factory：做出产品的地方）、\textbf{manufacture}（manu 手 + fact 做 $\to$ 手工制作）。"事实"就是"已经做成的事"。
\end{origin}

\vw{instance}{/ˈɪnstəns/}{n.}

\sense{n.}{例子，事例}
\col{an instance of sth}
\ex{This is a clear instance of discrimination.}
\col{in this/that instance 在这种情况下}
\ex{In this instance, we cannot accept the application.}
\col{in the first instance 首先，第一步}
\ex{In the first instance, contact your local office.}

\col{for instance 例如，比如}
\ex{Many countries, for instance France and Italy, produce excellent wines.}

\begin{origin}
来自拉丁语 \textit{instantia}（站到眼前、紧迫），源于 \textit{instare}（站在……之上、逼近）：\textit{in-}（在上）+ \textit{stare}（站立）。同源词：\textbf{instant}（瞬间，站在眼前的那一刻）、\textbf{insist}（坚持）、\textbf{stand}。"例子"与"瞬间"同源——都是"站到你眼前"的事物。
\end{origin}

\vw{differ}{/ˈdɪfə(r)/}{vi.}

\sense{vi.}{相异}
\col{differ from sth}
\ex{British English differs from American English in spelling.}
\col{differ in sth}
\ex{The twins differ widely in character.}

\sense{vi.}{意见相左}
\col{differ with/from sb on/about sth}
\ex{I differ with you on this point.}
\col{agree to differ 求同存异}
\ex{We'll never agree, so let's agree to differ.}
\col{I beg to differ 恕我不敢苟同（正式）}
\ex{I beg to differ — the evidence points the other way.}

\begin{origin}
来自拉丁语 \textit{differre}（搬开、使分开）：\textit{dis-}（分开）+ \textit{ferre}（携带）。同源词：\textbf{transfer}（trans- 越过 + fer 携带 $\to$ 转移）、\textbf{refer}、\textbf{offer}、\textbf{prefer}。\textit{ferre}（携带）是拉丁语中最多产的词根之一；"意见相左"就是"各执己见、各带一边"。
\end{origin}

\vw{conclusion}{/kənˈkluːʒn/}{n.}

\sense{n.}{结论，推论}
\col{come to/draw/reach a conclusion}
\ex{What conclusions can we draw from the data?}
\col{jump to conclusions 贸然下结论}
\ex{Don't jump to conclusions before you have all the facts.}

\sense{n.}{结束，结果}
\col{in conclusion 总之}
\ex{In conclusion, I would like to thank everyone involved.}
\col{bring sth to a conclusion 使……告一段落}
\ex{The meeting was brought to a conclusion at five o'clock.}

\sense{n.}{签订，达成}
\col{the conclusion of a treaty/agreement}
\ex{The conclusion of the trade agreement took months of negotiation.}

\begin{origin}
来自拉丁语 \textit{conclusio}，源于 \textit{concludere}（关闭、收拢）：\textit{con-}（一起）+ \textit{claudere}（关闭）。同源词：\textbf{include}（in 内 + claudere 关 $\to$ 关在里面 $\to$ 包含）、\textbf{exclude}（排除）、\textbf{close}。"把话收拢起来"就是结论。
\end{origin}

\vw{false}{/fɔːls/}{adj.}

\sense{adj.}{错误的，不真实的}
\col{false information/claims}
\ex{The report was based on false information.}
\col{give a false impression}
\ex{The media gave a false impression of the event.}

\sense{adj.}{非天生的}
\col{false teeth/hair/nails}
\ex{Grandpa wears false teeth.}

\sense{adj.}{假的，伪造的}
\col{a false passport/identity}
\ex{He travelled on a false passport.}
\col{a false alarm 虚惊一场}
\ex{The fire turned out to be a false alarm.}

\sense{adj.}{不真诚的}
\col{a false smile/promise}
\ex{She gave me a false smile and walked away.}
\col{under false pretences 以欺骗手段}
\ex{He obtained the money under false pretences.}

\begin{origin}
来自拉丁语 \textit{falsus}（欺骗的、假的），是 \textit{fallere}（欺骗）的过去分词。同源词：\textbf{fail}（failere 欺骗不成 $\to$ 失败）、\textbf{fault}（过错）、\textbf{fallacy}（谬误）、\textbf{fallible}（会犯错的）、\textbf{falsehood}（谎言）。"假"的词根是"骗"。
\end{origin}

\vp{bring sth to light}

\sense{v.}{揭露，披露}
\col{bring sth to light}
\ex{The investigation brought new evidence to light.}
\col{come to light 被揭露，为人所知}
\ex{More details have come to light since the trial.}
\col{shed/throw light on sth 阐明，使……清楚}
\ex{The documents shed new light on the affair.}

\begin{origin}
短语中的 light 指"公开、暴露"，源自《圣经》译本中"把暗中的事带到光下"的比喻（如《路加福音》"nothing is hidden that will not be brought to light"）。"曝光"古今中外同一种直觉。
\end{origin}

\vp{sum up}

\sense{v.}{总结，概括}
\col{to sum up 总之}
\ex{To sum up, the project has been a great success.}
\col{sum sth up}
\ex{He summed up the main points in a few sentences.}
\col{sum up the situation 迅速判断形势}
\ex{She quickly summed up the situation and took charge.}

\vw{sum}{/sʌm/}{n.}

\sense{n.}{金额，款项}
\col{a sum of money}
\ex{She inherited a large sum of money from her aunt.}
\col{a six-figure sum}
\ex{The painting sold for a six-figure sum.}

\sense{n.}{总和，总数}
\col{the sum of sth}
\ex{The sum of 2 and 3 is 5.}
\col{the sum total}
\ex{That's the sum total of my knowledge on the subject.}

\sense{n.}{全部}
\col{the sum of one's experience}
\ex{The sum of his life's experience is written in this book.}
\col{in sum 总之（书面）}
\ex{In sum, the advantages outweigh the disadvantages.}

\begin{origin}
来自拉丁语 \textit{summa}（最高处、顶端、总数），源于 \textit{summus}（最高的）。古罗马记账时把总额写在账目顶端，summa 遂由"顶端"转指"总数"。同源词：\textbf{summit}（顶峰）、\textbf{summary}（summary：把要点提到最上面 $\to$ 摘要）。"总数写在最高处"——一个关于记账习惯的词源。
\end{origin}

\vw{accurate}{/ˈækjərət/}{adj.}

\sense{adj.}{正确无误的}
\col{accurate information/figures}
\ex{The figures in the report are not accurate.}

\sense{adj.}{精确的}
\col{an accurate description/measurement}
\ex{The clock is accurate to within one second a year.}
\col{be accurate to within sth}
\ex{The device measures temperature accurate to within 0.1 degree.}

\begin{origin}
来自拉丁语 \textit{accuratus}（精心制作的），源于 \textit{accurare}（细心照料）：\textit{ad-}（去）+ \textit{cura}（照料）。同源词：\textbf{cure}（治愈）、\textbf{curious}、\textbf{secure}（se- 无 + cura 忧虑 $\to$ 无忧 $\to$ 安全）。"用心做出来的"自然"精确无误"。
\end{origin}

\vw{committed}{/kəˈmɪtɪd/}{adj.}

\sense{adj.}{尽心尽力的}
\col{a committed supporter/teacher}
\ex{She is a committed supporter of the environmental movement.}
\ex{The school employs committed and caring teachers.}

\col{be committed to sth/doing sth 致力于}
\ex{The company is committed to reducing its carbon footprint.}
\ex{We remain fully committed to the project.}

\begin{gram}[be committed to 中的 to 是介词]
\textit{be committed to} 里的 to 是介词，后面接名词或动名词，\textcolor{WildStrawberry!80!black}{\ding{55}} \textit{be committed to do sth}。同一规律的词组：\textit{be devoted to}（献身于）、\textit{be dedicated to}（致力于）、\textit{be addicted to}（沉迷于）、\textit{look forward to}（期待）。记忆口诀：这些 to 后面接"ing 的名词"。
\end{gram}

\vw{discrimination}{/dɪˌskrɪmɪˈneɪʃn/}{n.}

\sense{n.}{辨别力，识别力}
\col{discrimination between A and B}
\ex{He showed remarkable discrimination in choosing wines.}

\sense{n.}{歧视，区别对待}
\col{racial/sexual/age discrimination}
\ex{The law forbids discrimination on the grounds of race or gender.}
\col{discrimination against sb}
\ex{There is still discrimination against women in some industries.}
\col{discriminate (v.) 区别对待；歧视}
\ex{It is illegal to discriminate against older workers.}

\begin{origin}
来自拉丁语 \textit{discriminare}（分开），源于 \textit{discrimen}（间隔、区分），再追溯到 \textit{discernere}（辨别）：\textit{dis-}（分开）+ \textit{cernere}（筛选、区分）。同源词：\textbf{discern}（辨别）、\textbf{discernible}（可辨认的）。"辨别力"与"歧视"由同一个"区分"动作分化而来——区分得当是眼力，区分不公是歧视。
\end{origin}

\vw{curiosity}{/ˌkjʊəriˈɒsəti/}{n.}

\sense{n.}{好奇心，求知欲}
\col{curiosity about sth}
\ex{Children are full of curiosity about the world around them.}
\col{out of curiosity 出于好奇}
\ex{I opened the old box out of curiosity.}
\col{arouse/satisfy one's curiosity}
\ex{The letter aroused my curiosity.}
\col{Curiosity killed the cat. 好奇害死猫}

\sense{n.}{奇物，珍品}
\col{a curiosity shop 古玩店}
\ex{The old man ran a curiosity shop near the market.}
\ex{The exhibition displays curiosities from around the world.}

\begin{origin}
来自拉丁语 \textit{curiositas}（求知欲），源于 \textit{curiosus}（细心的、好打听的），词根 \textit{cura}（照料、操心）——与 \textbf{accurate} 同源。古义本是"对琐事过分操心"（含贬义，近"爱管闲事"），后来才转变成褒义的"好奇心"。
\end{origin}

\vw{commitment}{/kəˈmɪtmənt/}{n.}

\sense{n.}{献身，投入}
\col{commitment to sth/doing sth}
\ex{Her commitment to the cause of education is beyond doubt.}

\sense{n.}{承诺，许诺}
\col{make a commitment to do sth}
\ex{The government made a commitment to cut emissions by 2030.}
\col{honour/meet one's commitments 履行承诺}
\ex{The firm failed to meet its financial commitments.}

\sense{n.}{花费}
\col{a financial commitment 财政开支}
\ex{Buying a house is a major financial commitment.}
\col{a time commitment 时间投入}
\ex{The course requires a heavy time commitment.}

\begin{lex}{commitment / promise}
\textbf{promise} 是口头或书面承诺，重在"说"（\textit{make a promise}）；\textbf{commitment} 是更深的责任与投入，重在"做"（\textit{commitment to a cause} 对事业的献身）。商业与正式场合多用 commitment：\textit{meet one's commitments}（履行承诺）不能说 \textit{meet one's promises}。
\end{lex}

\vw{politics}{/ˈpɒlətɪks/}{n.}

\sense{n.}{政治，政治事务}
\col{enter/go into politics 从政}
\ex{She entered politics at the age of thirty.}
\col{local/national politics}
\ex{He has been involved in local politics for years.}
\col{office politics 办公室政治}
\ex{I'm tired of office politics.}

\sense{n.}{（个人的）政治观点，政见}
\col{one's politics}
\ex{His politics are rather left-wing.}

\begin{gram}[politics 的单复数谓语]
\textit{politics} 形式上以 -s 结尾，但指"政治、政治学、政治事务"时谓语用单数：\textit{Politics is a complicated business.}（政治是件复杂的事）。指"（个人的）政见"时谓语反而用复数：\textit{His politics are conservative.} 同类词：\textbf{mathematics}、\textbf{physics}、\textbf{economics}——学科名单数，指行为或观点时复数。
\end{gram}

\begin{origin}
来自希腊语 \textit{ta politika}（城邦事务），源于 \textit{polites}（公民），词根 \textit{polis}（城邦）。同源词：\textbf{policy}（政策）、\textbf{police}（police：维持城邦秩序的人）、\textbf{metropolis}（metro 母亲 + polis 城 $\to$ 母城 $\to$ 大都市）、\textbf{cosmopolitan}（cosmos 宇宙 + polites 公民 $\to$ 世界公民的）。亚里士多德那部《政治学》书名就叫 \textit{Politika}。
\end{origin}

\vw{accuse}{/əˈkjuːz/}{vt.}

\sense{vt.}{谴责，控诉}
\col{accuse sb of (doing) sth}
\ex{He was accused of stealing company secrets.}
\col{accuse sb of lying}
\ex{She accused him of lying to her.}
\col{the accused 被告}
\ex{The accused pleaded not guilty.}

\begin{lex}{accuse / charge / blame}
\textbf{accuse sb of sth}（控告、指责，最正式）：\textit{accuse him of theft}；\textbf{charge sb with sth}（法律起诉）：\textit{charge him with murder}；\textbf{blame sb for sth}（日常责怪）：\textit{blame me for the mistake}。介词分别为 of / with / for，是考试常考点。
\end{lex}

\begin{origin}
来自拉丁语 \textit{accusare}（传讯、起诉）：\textit{ad-}（去）+ \textit{causa}（讼案、原因）。同源词：\textbf{cause}（原因）、\textbf{excuse}（ex- 开脱 + causa 讼案 $\to$ 免于追责 $\to$ 借口）。"把某人告上法庭"就是 accuse 的本义。
\end{origin}

\vw{tax}{/tæks/}{n. \& vt.}

\sense{n.}{税，税款}
\col{income/property/sales tax}
\ex{The government raised income tax last year.}
\col{before/after tax}
\ex{She earns \pounds 40,000 a year before tax.}
\ex{Goods bought online are taxed at 20\%.}
\col{tax cuts/increases}
\ex{The new policy promises tax cuts for small businesses.}

\sense{vt.}{对……征税，使纳税}
\ex{The rich should be taxed more heavily.}

\col{tax sb's patience/stamina 考验某人的耐心/体力}
\ex{The endless questions taxed his patience.}

\begin{origin}
来自拉丁语 \textit{taxare}（触摸、估价、责备）——这是 \textit{tangere}（触摸）的反复动词（反复体）。同源词：\textbf{touch}（触摸）、\textbf{contact}（con- 一起 + tact 触 $\to$ 接触）、\textbf{tactile}（触觉的）。"摸一摸估个价"——收税最初就是摸摸财产估个值。
\end{origin}

\vw{mount}{/maʊnt/}{vi. \& vt. \& n.}

\sense{vi.}{逐步增加}
\col{mounting pressure/evidence}
\ex{Pressure is mounting on the government to act.}
\col{tension/doubts mount}
\ex{As the deadline approached, doubts mounted.}

\sense{vt.}{攀登，爬上}
\col{mount a horse/bicycle}
\ex{She mounted her horse and rode away.}
\col{mount the stairs}
\ex{He mounted the stairs slowly.}

\sense{n.}{山}
\col{Mount/Mt. + 山名}
\ex{Mount Qomolangma is the highest peak on Earth.}

\col{mount an exhibition/campaign 举办展览；发动运动}
\ex{The museum mounted an exhibition of modern art.}

\col{mount up（尺寸和数量上）增加，上升}
\ex{The bills are mounting up.}
\ex{Evidence against him mounted up quickly.}

\begin{origin}
名词"山"与动词"登上"本是一家人：拉丁语 \textit{mons}（山）派生出动词 \textit{montare}（登上）。同源词：\textbf{mountain}（山）、\textbf{amount}（a- 向 + mont 山顶 $\to$ 累积到顶点 $\to$ 总额）、\textbf{surmount}（sur- 越过 + mont $\to$ 克服、登上）、\textbf{paramount}（par- 通过 + amont $\to$ 至高的）。
\end{origin}

\vw{elevation}{/ˌelɪˈveɪʃn/}{n.}

\sense{n.}{提高；提升}
\col{the elevation of sb to a post 晋升某人到某职}
\ex{His elevation to the board came as a surprise.}

\sense{n.}{海拔}
\col{at an elevation of ...}
\ex{The city lies at an elevation of 1,200 metres.}

\begin{origin}
来自拉丁语 \textit{elevare}（举起）：\textit{e-}（向外、向上）+ \textit{levare}（使变轻、举起），词根 \textit{levis}（轻的）。同源词：\textbf{elevator}（美式英语的电梯）、\textbf{lever}（杠杆）、\textbf{levity}（轻浮）、\textbf{relieve}（减轻）。"变轻"才"举得起"——古人用同一词根表达"轻"与"升"。
\end{origin}

\vw{profession}{/prəˈfeʃn/}{n.}

\sense{n.}{行业，职业}
\col{the teaching/medical/legal profession}
\ex{She is a lawyer by profession.}
\col{enter/go into a profession}
\ex{He entered the teaching profession in 2010.}

\sense{n.}{同行}
\col{the profession（统称）本行业全体人员}
\ex{The profession as a whole opposes the new rules.}

\sense{n.}{宣称，声明}
\col{a profession of faith/loyalty 信仰表白；效忠声明}
\ex{He made a profession of loyalty to the king.}

\begin{lex}{profession / occupation / career / job}
\textbf{profession} 指须经专门训练、常需执照的专业职业（律师、医生、教师）；\textbf{occupation} 泛指任何职业，是表格用词（\textit{What's your occupation?}）；\textbf{career} 指职业生涯、事业轨迹；\textbf{job} 是具体的一份工作。
\end{lex}

\begin{origin}
来自拉丁语 \textit{professio}（公开声明），源于 \textit{profiteri}（当众宣布）：\textit{pro-}（当众、向前）+ \textit{fateri}（承认、说出）。同源词：\textbf{professor}（professor：当众宣讲的人）、\textbf{confess}（con- 完全 + fateri $\to$ 彻底说出 $\to$ 忏悔、承认）。中世纪大学里，入行前须当众宣誓，profession 遂由"公开宣誓"转指"专业职业"。
\end{origin}

\vw{drama}{/ˈdrɑːmə/}{n.}

\sense{n.}{戏剧}
\col{a TV/radio drama}
\ex{She studied drama at college.}
\col{a costume drama 古装剧}
\ex{The costume drama attracted millions of viewers.}

\sense{n.}{戏剧性事件}
\col{high drama 高潮迭起的场面}
\ex{The election campaign was full of high drama.}
\col{a drama queen 小题大做的人（口语）}
\ex{Stop being such a drama queen!}

\begin{origin}
来自希腊语 \textit{drama}（行动、戏剧），源于 \textit{dran}（做、行动）。亚里士多德在《诗学》中指出，悲剧是对"行动"（praxis）的模仿——戏剧就是把"做出来的事"搬上舞台。同源词：\textbf{dramatic}（戏剧性的）。
\end{origin}

\vw{nevertheless}{/ˌnevəðəˈles/}{adv.}

\sense{adv.}{尽管如此}
\col{句首用法（后接逗号）}
\ex{The weather was terrible. Nevertheless, the match went ahead.}
\col{句中、句尾用法}
\ex{It was hard work. We enjoyed it nevertheless.}
\ex{The book is old; it is nevertheless worth reading.}

\begin{lex}{nevertheless / nonetheless / however}
\textbf{nevertheless} 与 \textbf{nonetheless} 同义且都可互换，书面色彩浓，二者同源于 never the less（一点都不少）；\textbf{however} 更常用、更口语化，三词都表"尽管如此、然而"。注意 however 也可作连接副词放句首。
\end{lex}

\begin{origin}
中古英语 \textit{never the less}（never + the + less）字面义是"一点也不少"——本义指"数量上毫无减损"，后来反讽式地转义为"即使如此"，成为让步副词。与 nonetheless 完全同构，是典型的"短语凝固成词"。
\end{origin}

\vw{witness}{/ˈwɪtnəs/}{vt. \& vi. \& n.}

\sense{vt.}{目击}
\col{witness a crime/accident}
\ex{Several people witnessed the robbery.}

\sense{vt.}{见证；是发生……的时间（或地点）}
\col{sth witness sth（无生命主语）}
\ex{The stadium has witnessed countless victories.}
\ex{The past decade witnessed rapid changes in technology.}

\sense{vt.}{作证}
\col{witness against/for sb}
\ex{He refused to witness against his own brother.}

\sense{vt.}{是……的迹象}
\ex{The empty streets witnessed the severity of the lockdown.}

\sense{n.}{目击者，见证人；证人}
\col{a witness to/of sth}
\ex{She was the only witness to the accident.}
\col{bear witness to sth 证明，为……作证}
\ex{His achievements bear witness to his talent.}
\col{a key witness 关键证人}
\ex{The key witness gave evidence in court.}

\begin{gram}[witness 的无生命主语用法]
\textit{witness} 常用于"时间/地点作主语"的句式：\textit{The 20th century witnessed two world wars.} 这类"无生命主语 + 见证类动词"的拟人化结构是英语书面语一大特色，同类动词还有 \textbf{see}、\textbf{find}：\textit{The city saw a rapid recovery.} 汉语更习惯用"在……发生了"，翻译时要转换视角。
\end{gram}

\begin{origin}
源自古英语 \textit{witness}（证词、证据）：\textit{wit}（知识、心智，动词 \textit{witan} 知道）+ 名词后缀 -ness。同源词：\textbf{wit}（机智）、\textbf{wise}（明智）、\textbf{wisdom}（智慧）。词根可追溯到原始印欧语 *\textit{weid-}（看见、知道）——与拉丁语 \textit{videre}（看见，即 \textbf{video} 的词根）同源。证人，就是"亲眼所见、知道真相的人"。
\end{origin}

\vw{edition}{/ɪˈdɪʃn/}{n.}

\sense{n.}{版本}
\col{a new/revised edition}
\ex{The new edition of the dictionary has 5,000 more entries.}
\col{a limited edition 限量版}
\ex{The book was published in a limited edition of 1,000 copies.}

\sense{n.}{一份，一期}
\col{the morning/evening edition}
\ex{I read the morning edition of the paper.}

\sense{n.}{版次}
\col{a first edition 初版}
\ex{Collectors pay a fortune for a first edition.}

\begin{origin}
来自拉丁语 \textit{editio}（出版、发行），源于 \textit{edere}（发出、出版）：\textit{e-}（向外）+ \textit{dare}（给）。同源词：\textbf{editor}（编辑）、\textbf{edit}（编辑）。"出版"字面义即"把书交给公众"。
\end{origin}

\vw{platform}{/ˈplætfɔːm/}{n.}

\sense{n.}{计算机平台}
\col{a digital/social media platform}
\ex{The app works on both major platforms.}

\sense{n.}{站台，月台}
\col{platform + 数字}
\ex{The train to Paris leaves from platform 5.}

\sense{n.}{讲台，舞台}
\col{speak/mount from a platform}
\ex{The speaker mounted the platform and began his speech.}

\col{a political platform 政纲，竞选纲领}
\ex{The party's election platform promises tax cuts.}

\begin{origin}
16 世纪借自法语 \textit{plateforme}（扁平的形状）：\textit{plate}（扁平）+ \textit{forme}（形状）。plate 最终源自希腊语 \textit{platys}（宽阔的、扁平的）——同源词：\textbf{plate}（盘子）、\textbf{plaza}（广场）。平台就是"一块平的东西"。
\end{origin}

\vp{spring up}

\sense{v.}{迅速出现，突然兴起}
\ex{New shopping centres have sprung up all over the city.}
\ex{Towns sprang up along the railway line.}
\ex{Suspicion sprang up in his mind.}

\begin{origin}
spring 本义"跳跃、迸出"（古英语 \textit{springan}，同源德语 \textit{springen} 跳），引申出"泉水涌出、萌芽生长"。同源词：\textbf{spring}（泉水、弹簧、春天——万物迸发的季节）。spring up 即"像泉水一样一下子冒出来"。
\end{origin}

\vw{interaction}{/ˌɪntərˈækʃn/}{n.}

\sense{n.}{互动}
\col{interaction between/with}
\ex{Social interaction is important for children's development.}
\ex{There is little interaction between the two communities.}

\sense{n.}{互相作用}
\col{the interaction of/among}
\ex{The disease is caused by the interaction of genes and environment.}

\begin{origin}
来自拉丁语：\textit{inter}（在……之间）+ \textit{actio}（行动），后者源于 \textit{agere}（做、驱动）。同源词：\textbf{action}、\textbf{active}、\textbf{agent}（代理人）、\textbf{actor}。"彼此之间的作用"即互动。
\end{origin}

\vw{chart}{/tʃɑːt/}{n. \& vt.}

\sense{n.}{图表}
\col{a bar/pie/line chart}
\ex{The chart shows the growth of online sales.}

\sense{n.}{海图}
\col{a nautical chart 航海图}
\ex{The captain studied the nautical chart carefully.}

\sense{n.}{(the charts) 排行榜}
\col{top the charts 登上排行榜首位}
\ex{The song topped the charts for six weeks.}

\sense{vt.}{记录}
\col{chart the progress/development of}
\ex{The book charts the rise and fall of the empire.}

\sense{vt.}{制订计划}
\col{chart a course}
\ex{We charted a course for the future.}

\sense{vt.}{绘制区域地图}
\col{chart the coastline}
\ex{The expedition charted the unknown coastline.}

\begin{origin}
来自拉丁语 \textit{charta}（纸、莎草纸），再追溯到希腊语 \textit{kharēs}（莎草叶、纸张）。同源词：\textbf{card}（卡片）、\textbf{cartoon}（cartoon：画在纸上的草图）、\textbf{charter}（特许状）、\textbf{cartography}（制图学）。"图表"最初就是"画在纸上"。
\end{origin}

\vw{broadcast}{/ˈbrɔːdkɑːst/}{n. \& vi. \& vt.}

\sense{n.}{广播节目，电视节目}
\col{a live broadcast}
\ex{Millions watched the live broadcast of the final.}

\sense{vt. \& vi.}{播送，广播}
\col{broadcast (broadcast, broadcast)}
\ex{The match was broadcast to over 100 countries.}
\ex{The station broadcasts news every hour.}

\sense{vt.}{散布，传播}
\col{broadcast a message/rumour}
\ex{Don't broadcast your plans to everyone.}

\begin{gram}[不变形动词：broadcast–broadcast–broadcast]
\textit{broadcast} 的过去式和过去分词与动词原形相同（\textit{It was broadcast}，而非 broadcasted）。同类"三式相同"的动词：\textbf{cost}、\textbf{cut}、\textbf{hit}、\textbf{hurt}、\textbf{let}、\textbf{put}、\textbf{set}、\textbf{shut}。这些多为单音节短动词，broadcast 是其中少见的复合词。
\end{gram}

\begin{origin}
由 broad（宽广的）+ cast（抛撒）构成。cast 本义"扔、撒"（借自古诺斯语 \textit{kasta} 抛掷）。农业上有 broadcast seeding（撒播），指把种子撒向大片田地；20 世纪无线电兴起后，"把信号撒向四方"便成了广播。同源词：\textbf{forecast}（预先抛出的预测）、\textbf{cast}。
\end{origin}

\vw{persuasion}{/pəˈsweɪʒn/}{n.}

\sense{n.}{说服，劝说}
\col{the power/art of persuasion}
\ex{She has the power of persuasion.}
\col{after much persuasion}
\ex{After a lot of persuasion, she agreed to come.}

\sense{n.}{信仰}
\col{of different persuasions}
\ex{Voters of all political persuasions support the plan.}

\begin{origin}
来自拉丁语 \textit{persuasio}，源于 \textit{persuadere}（彻底说动）：\textit{per-}（彻底）+ \textit{suadere}（劝告）。词根可追溯到原始印欧语 *\textit{swād-}（甜）——与希腊语 \textit{hedys}（甜）同源（\textbf{hedonism} 享乐主义）。"说服"就是用甜言蜜语把人说动。
\end{origin}

\vw{persuade}{/pəˈsweɪd/}{vt.}

\sense{vt.}{说服}
\col{persuade sb to do sth}
\ex{She persuaded her father to buy the house.}
\col{persuade sb into/out of doing sth}
\ex{He persuaded me out of quitting my job.}

\sense{vt.}{使信服}
\col{persuade sb of sth}
\ex{They persuaded me of the truth of the story.}
\col{persuade sb that ...}
\ex{The doctor persuaded me that the operation was necessary.}

\begin{lex}{persuade / convince}
\textbf{persuade} 侧重"说服某人采取行动"（\textit{persuade sb to do sth}）；\textbf{convince} 侧重"使某人相信"（\textit{convince sb that ...} / \textit{convince sb of sth}）。但美式英语里 convince 也可接不定式：\textit{He convinced me to stay.} 英式英语考试中仍倾向区分二者。
\end{lex}

\vw{advertisement}{/ədˈvɜːtɪsmənt/}{n.}

\sense{n.}{广告，启事}
\col{put/place an advertisement}
\ex{She placed an advertisement in the local newspaper.}
\col{a job advertisement}
\ex{They are running a job advertisement online.}
\col{(informal) ad 广告}
\ex{an ad for a new car}

\sense{n.}{广告活动，广告宣传}
\col{an advertisement for sth}
\ex{The commercial was an advertisement for a new soft drink.}

\begin{origin}
advertise 的古法语来源 \textit{avertir}（提醒、引人注意）最终可追溯至拉丁语 \textit{advertere}（转向）：\textit{ad-}（朝向）+ \textit{vertere}（转）。同源词：\textbf{convert}（转换）、\textbf{reverse}（反转）、\textbf{verse}（诗行，转回下一行）。广告就是把"注意力转过来"。
\end{origin}

\vw{channel}{/ˈtʃænl/}{n.}

\sense{n.}{途径，渠道}
\col{go through the proper channels 走正当渠道}
\ex{We must go through the proper channels.}
\col{official/diplomatic channels}
\ex{The two countries communicated through diplomatic channels.}
\col{a distribution channel 分销渠道}
\ex{Online stores are now a major distribution channel.}

\sense{n.}{电视台；频道}
\col{switch/change channels}
\ex{He kept switching channels during the ad break.}

\sense{n.}{方法}
\ex{Email is a quick channel for sending messages.}

\sense{n.}{水渠}
\col{an irrigation channel 灌溉渠}
\ex{The fields are watered through irrigation channels.}

\col{the English Channel 英吉利海峡}
\ex{They swam across the English Channel.}

\begin{origin}
来自拉丁语 \textit{canalis}（管道、沟渠），源于 \textit{canna}（芦苇）——再往前是希腊语 \textit{kanna}（芦苇）。同源词：\textbf{canal}（运河）、\textbf{cane}（藤条、甘蔗）。芦苇茎中空可通水，于是"芦苇管"成了"管道"，又引申为传递信息的"渠道、频道"。
\end{origin}

\vw{peak}{/piːk/}{n. \& adj. \& vi.}

\sense{n.}{顶峰，高峰}
\col{at the peak of sth}
\ex{She was at the peak of her career.}
\col{reach a peak}
\ex{Sales reached a peak in December.}

\sense{n.}{山峰}
\ex{The mountain peak was covered with snow.}

\sense{n.}{尖端}
\col{the peak of a roof}
\ex{Birds perched on the peak of the roof.}

\sense{adj.}{高峰时期的，最高度的}
\col{peak hours/season 高峰时段；旺季}
\ex{Traffic is worst during peak hours.}
\col{peak demand}
\ex{Electricity prices rise at times of peak demand.}

\sense{vi.}{达到高峰，达到最高值}
\col{peak at sth}
\ex{The temperature peaked at 40 degrees.}
\ex{His popularity peaked in the 1990s.}

\begin{origin}
很可能来自 pike（尖物）的变体——古英语 \textit{pīc}（尖头物）。同源词：\textbf{pike}（长矛、梭鱼）、\textbf{pick}（镐）。"像矛尖一样的山顶"就是 peak，进而引申出"最高点、高峰期"。
\end{origin}

\vw{advertise}{/ˈædvətaɪz/}{vt. \& vi.}

\sense{vt. \& vi.}{做广告，登广告}
\col{advertise sth}
\ex{The company advertises its products on television.}
\col{advertise heavily/widely}
\ex{The new phone was advertised heavily before launch.}

\sense{vt.}{公布，征聘}
\col{advertise for sth}
\ex{We advertised for a new assistant.}

\sense{vt.}{展现，宣传}
\col{advertise oneself as sth}
\ex{The hotel advertises itself as eco-friendly.}

\begin{origin}
与 advertisement 同源：拉丁语 \textit{advertere}（转向）$\to$ 古法语 \textit{avertir}（引人注意）$\to$ 英语 advertise。"做广告"即"把目光吸引过来"。
\end{origin}

\vw{boost}{/buːst/}{vt. \& n.}

\sense{vt.}{使增长，使兴旺}
\col{boost sales/production/the economy}
\ex{The advertising campaign boosted sales by 20\%.}
\col{boost one's confidence/morale}
\ex{The early goal boosted the team's confidence.}

\sense{n.}{增长，提高}
\col{a boost in sth}
\ex{We saw a boost in profits last quarter.}

\sense{n.}{帮助，激励}
\col{a boost to/for sth}
\ex{The victory was a huge boost to the team.}
\col{give sb a boost}
\ex{Winning the award gave her career a real boost.}

\begin{origin}
词源并不十分确定：19 世纪初出现在美式英语中，一说与方言词 boom（突然增长、兴旺）有关。词源不明的词反而常是高频词——英语里不少日常动词（如 \textbf{big}、\textbf{boy}）都查不到确切祖先。
\end{origin}

\vw{purchase}{/ˈpɜːtʃəs/}{vt. \& n.}

\sense{vt.}{买，购买}
\col{purchase sth from sb/sth（正式）}
\ex{The firm purchased a new computer system from a local supplier.}

\sense{n.}{购买，采购}
\col{the purchase of sth}
\ex{the purchase of a new car}

\sense{n.}{购买的东西，购买项目}
\col{make a purchase}
\ex{She made several purchases online.}
\col{a major/large purchase}
\ex{A car is a major purchase.}

\col{purchasing power 购买力}
\ex{Rising prices have weakened consumers' purchasing power.}

\begin{origin}
来自古法语 \textit{purchacier}（设法弄到手）：\textit{pur-}（为了）+ \textit{chacier}（追逐）——chacier 正是现代英语 \textbf{chase}（追逐）的来源。同源词：\textbf{chase}。"购买"字面义是"追着要弄到手"。口语中"买"用 buy，purchase 偏正式。
\end{origin}

\vp{get across}

\sense{v.}{被传达，被理解；把……讲清楚}
\col{get sth across to sb}
\ex{He found it hard to get his ideas across to the audience.}
\col{get across（不及物）}
\ex{The message didn't get across.}
\ex{Teachers must get their points across clearly.}

\col{同义：put sth across（表达清楚）；communicate sth}

\vw{brand}{/brænd/}{n.}

\sense{n.}{品牌}
\col{a brand name}
\ex{Which brand of coffee do you prefer?}
\col{brand loyalty/awareness 品牌忠诚度；品牌知名度}
\ex{Strong brand loyalty is hard to break.}

\sense{n.}{类型}
\col{a brand of sth}
\ex{He has his own brand of humour.}

\sense{n.}{烙印}
\col{a cattle brand}
\ex{Cattle are marked with the ranch's brand.}

\col{brand new 崭新的}
\ex{He arrived in a brand new sports car.}

\begin{origin}
源自古英语 \textit{brand}（燃烧的木块、火炬），词根为日耳曼语"烧"。同源词：\textbf{burn}（燃烧）、\textbf{brandy}（白兰地——把葡萄酒"烧"过蒸馏而成）、\textbf{firebrand}（火把）。牧场主用烧红的烙铁给牲口打"烙印"，19 世纪商标兴起后，brand 遂由"烙印"转指"品牌"。
\end{origin}

\vw{ambassador}{/æmˈbæsədə(r)/}{n.}

\sense{n.}{大使，代表}
\col{ambassador to + 国家}
\ex{She was appointed ambassador to China.}
\col{an ambassador of/for sth}
\ex{He serves as an ambassador for world peace.}
\col{a brand ambassador 品牌代言人}
\ex{The singer became the brand ambassador of a luxury watchmaker.}

\begin{origin}
经由法语 \textit{ambassadeur} 进入英语，再往前是意大利语 \textit{ambasciatore}，最终源自高卢语（凯尔特语）\textit{ambactus}（仆人、随从），经日耳曼语中转。同源词：\textbf{embassy}（大使馆）。"被差遣的仆人"——使节最初不过是"跑腿送信的人"。
\end{origin}

\vw{rating}{/ˈreɪtɪŋ/}{n.}

\sense{n.}{收视率}
\col{the ratings（复数）收视率}
\ex{The show has high ratings.}

\sense{n.}{等级，级别}
\col{a rating system}
\ex{The film has an age rating of 12.}
\col{a credit rating 信用评级}
\ex{The country's credit rating was downgraded.}
\col{approval ratings 支持率}
\ex{The president's approval ratings fell sharply.}

\begin{origin}
rate 来自古法语 \textit{rate}（价格、估值），源于拉丁语 \textit{rata} (pars)（固定的份额），词根 \textit{reri}（思考、估定）。同源词：\textbf{ratio}（比率）、\textbf{rational}（理性的）、\textbf{reason}（reason：估定之事 $\to$ 理由）。"评级"就是"估定一个档次"。
\end{origin}

\vw{sponsor}{/ˈspɒnsə(r)/}{vt. \& n.}

\sense{vt.}{赞助}
\col{sponsor a team/event}
\ex{The bank sponsors the local football team.}

\sense{vt.}{主办}
\col{sponsor a conference}
\ex{The university sponsored an international conference on AI.}

\sense{vt.}{为慈善活动捐款}
\col{sponsor sb for a charity}
\ex{She is sponsoring a colleague for a charity run.}

\sense{vt.}{资助}
\col{sponsor a student}
\ex{He sponsored a student through medical school.}

\sense{n.}{赞助商；为慈善活动捐款的人}
\col{the sponsor of sth}
\ex{The main sponsor of the marathon is a sports brand.}
\col{corporate sponsors 企业赞助商}
\ex{Corporate sponsors provide most of the funding.}

\begin{origin}
来自拉丁语 \textit{spondere}（郑重许诺、立誓担保），sponsor 即"担保人"。同源词：\textbf{respond}（respond：re- 回 + spondere 许诺 $\to$ 应诺 $\to$ 回应）、\textbf{response}（回应）、\textbf{spouse}（配偶，立过婚誓的人）、\textbf{espouse}（支持）。"出钱担保承诺的人"就是赞助人。
\end{origin}

\vw{absorb}{/əbˈzɔːb/}{vt.}

\sense{vt.}{理解，掌握}
\col{absorb knowledge/information}
\ex{Children absorb new languages very quickly.}
\ex{She absorbed every detail of the lecture.}

\sense{vt.}{吸收}
\col{absorb water/heat/light}
\ex{Plants absorb carbon dioxide from the air.}
\col{be absorbed into 被并入}
\ex{The small company was absorbed into a large corporation.}

\sense{vt.}{吸引全部注意力}
\col{be absorbed in sth 全神贯注于}
\ex{He was completely absorbed in his book.}

\begin{origin}
来自拉丁语 \textit{absorbere}（吞没、吸进）：\textit{ab-}（离开、去）+ \textit{sorbere}（吸、啜）。同源词：\textbf{sorbet}（果汁冰糕——吸着喝的甜食）、\textbf{absorbent}（吸水的）。"吸进去"便同时有了"吸收营养"和"吸收知识"两个方向。
\end{origin}

\vw{tailor}{/ˈteɪlə(r)/}{vt. \& n.}

\sense{vt.}{专门制作，定做}
\col{tailor sth to sth/sb}
\ex{The course is tailored to the needs of beginners.}
\col{tailor-made 定制的；恰到好处的}
\ex{We offer tailor-made training programmes.}

\sense{n.}{裁缝}
\col{a tailor's 裁缝店}
\ex{He had the suit made at a tailor's.}

\begin{origin}
来自古法语 \textit{tailleur}（裁剪者），源于 \textit{taillier}（剪裁），再追溯到拉丁语 \textit{taliare}（切割），词根 \textit{talea}（枝条、切块）。同源词：\textbf{detail}（de- 切碎 + tailler $\to$ 切成小块 $\to$ 细节）、\textbf{retail}（re- 再 + tailler $\to$ 再切开 $\to$ 零售）。retail 字面义就是"把整批货切开来零卖"——与裁缝是同一家人。
\end{origin}

\vp{brighten up}

\sense{v.}{（使）增添乐趣；使更艳丽}
\col{brighten (sth) up}
\ex{Flowers brighten up any room.}
\ex{Her smile brightened up the whole meeting.}
\col{brighten up（天气）放晴}
\ex{The sky is brightening up.}

\begin{origin}
bright 源自古英语 \textit{beorht}（明亮的），加使动后缀 -en（让……变……）构成动词。同类后缀词：\textbf{darken}、\textbf{deepen}、\textbf{widen}、\textbf{sharpen}、\textbf{strengthen}。brighten up 即"让……亮起来"，引申为"添彩、提神"。
\end{origin}

\vw{housing}{/ˈhaʊzɪŋ/}{n.}

\sense{n.}{住房，住宅}
\col{the housing market}
\ex{Housing prices have risen sharply.}
\col{affordable housing 经济适用住房}
\ex{The city is short of affordable housing.}
\col{a housing shortage 住房短缺}
\ex{The housing shortage affects young people most.}

\begin{origin}
house 源自古英语 \textit{hūs}，可能源自原始印欧语 *(s)\textit{keu-}（遮盖、隐藏）——同源词：\textbf{hide}（藏）、\textbf{hut}（小屋）。房子最初就是"遮风蔽日之所"，housing 由"房屋"统称引申为"住房供给"。
\end{origin}

\vw{estate}{/ɪˈsteɪt/}{n.}

\sense{n.}{住宅区}
\col{a housing estate 住宅区，住宅群}
\ex{They live on a new housing estate.}

\sense{n.}{（通常指农村的）大片私有土地}
\col{a country/large estate}
\ex{He inherited a country estate with 200 acres.}

\sense{n.}{个人财产，（尤指）遗产}
\col{real estate 房地产}
\ex{Real estate prices in the city have doubled.}
\col{an estate agent 房产中介}
\ex{The estate agent showed us three flats.}

\begin{origin}
来自拉丁语 \textit{status}（状态、状况）$\to$ 古法语 \textit{estat}（状况、身份）$\to$ 英语 estate。它和 \textbf{state} 是同一拉丁词源分出的"双重词"（doublet）。"财产状况"逐渐缩指"地产、庄园"——尤指作为财富的地产。
\end{origin}

\vw{amuse}{/əˈmjuːz/}{vt.}

\sense{vt.}{（提供）消遣}
\col{amuse oneself with sth}
\ex{He amused himself by playing computer games.}

\sense{vt.}{逗笑}
\col{amuse sb}
\ex{The clown amused the children for hours.}
\col{be amused at/by sth}
\ex{I was amused by his little jokes.}

\col{amusing (adj.) 有趣的}
\ex{an amusing story}

\begin{origin}
来自古法语 \textit{amuser}（使发呆、消磨时间）：\textit{a-}（使）+ \textit{muser}（凝视、沉思）。同源词：\textbf{muse}（缪斯女神；沉思）、\textbf{museum}（缪斯的神殿）。"让人看得出神"就是消遣、逗乐——amuse 与 museum 竟是同根。
\end{origin}

% ---------------- 选收的未标 * 词条 ----------------

\vw{dozen}{/ˈdʌzn/}{n.}

\sense{n.}{（一）打，十二个；十来个}
\col{a dozen 一打}
\ex{a dozen eggs}
\col{half a dozen 半打}
\ex{I'll take half a dozen, please.}
\col{dozens of 很多，许多}
\ex{There were dozens of people at the party.}
\col{talk nineteen to the dozen 喋喋不休}
\ex{The two friends talked nineteen to the dozen all evening.}

\begin{origin}
来自古法语 \textit{dozeine}（十二个），源于拉丁语 \textit{duodecim}（十二）：\textit{duo}（二）+ \textit{decem}（十）。同源词：\textbf{duodecimal}（十二进制的）。"一打"是诺曼人带给英语的计数单位。
\end{origin}

\vw{found}{/faʊnd/}{vt.}

\sense{vt.}{创建，创办；建立，兴建}
\col{found a company/school}
\ex{The university was founded in 1895.}
\col{founder (n.) 创始人}
\ex{the founder of the company}

\col{注意：find–found–found（发现）与 found–founded–founded（创建）拼写相同，词源无关}

\begin{origin}
来自拉丁语 \textit{fundare}（打地基），词根 \textit{fundus}（底部）。同源词：\textbf{fund}（基金）、\textbf{foundation}（地基、基金会）、\textbf{profound}（pro- 前 + fundus 底 $\to$ 深及底部的 $\to$ 深刻的）。"打下地基"就是创立。
\end{origin}

\vw{scholarship}{/ˈskɒləʃɪp/}{n.}

\sense{n.}{奖学金}
\col{win/award a scholarship}
\ex{She won a scholarship to study at Cambridge.}

\sense{n.}{学术}
\ex{His scholarship is beyond question.}

\begin{origin}
scholar（学者）来自拉丁语 \textit{schola}（学校），再追溯到希腊语 \textit{skholē}——本义竟是"闲暇"！古希腊人认为"有空闲"才能读书做学问。同源词：\textbf{school}、\textbf{scholar}。-ship 是表"身份、状态"的名词后缀（同 friendship、leadership）。
\end{origin}

\vw{category}{/ˈkætəɡəri/}{n.}

\sense{n.}{类别，种类}
\col{fall into a category}
\ex{These crimes fall into a separate category.}
\col{in a category of its own 独树一帜}
\ex{This wine is in a category of its own.}

\begin{origin}
来自希腊语 \textit{kategoria}，本义是"公开指控、断言"，由 \textit{kata}（向下）+ \textit{agoreuein}（在广场上讲话）构成，\textit{agora} 即"集市"。亚里士多德《范畴篇》把 kategoria 用作逻辑术语，赋予其"类别"义——"在集市上公开断言某物为何物"，就是给它分类。
\end{origin}

\vw{psychology}{/saɪˈkɒlədʒi/}{n.}

\sense{n.}{心理学}
\ex{She studied psychology at university.}

\sense{n.}{心理，心理特征}
\col{the psychology of sth}
\ex{the psychology of crowds}

\begin{origin}
来自希腊语 \textit{psukhē}（灵魂、呼吸）+ \textit{-logia}（学问）。\textit{psukhē} 本义是"呼吸"——生命的气息。同源词：\textbf{psychiatrist}（精神科医生）、\textbf{psychic}（通灵的）。从"一口气"到"灵魂"再到"心理"，一词跨越三千年。
\end{origin}

\vw{memorable}{/ˈmemərəbl/}{adj.}

\sense{adj.}{难忘的，值得纪念的}
\col{a memorable experience/occasion}
\ex{It was a memorable evening.}

\begin{origin}
来自拉丁语 \textit{memorabilis}（值得记住的），源于 \textit{memorare}（记得），词根 \textit{memor}（记在心里的）。同源词：\textbf{memory}（记忆）、\textbf{remember}（remember：re- 再 + memor 记 $\to$ 重新记起）、\textbf{memorial}（纪念碑）。
\end{origin}

\vw{slogan}{/ˈsləʊɡən/}{n.}

\sense{n.}{标语，口号}
\col{a campaign slogan 竞选口号}
\ex{The party's slogan was "Change now".}

\begin{origin}
来自苏格兰盖尔语 \textit{sluagh-ghairm}（战斗呐喊）：\textit{sluagh}（军队）+ \textit{gairm}（呼喊）。古代高地部族开战前齐声高喊部族口号，slogan 由此而来——今天广告人的标语，原是战场上的战吼。
\end{origin}

\vw{discount}{/ˈdɪskaʊnt/}{n. \& vt.}

\sense{n.}{折扣}
\col{at a discount 打折}
\ex{The store sells last season's clothes at a discount.}
\col{a discount on sth}
\ex{Members get a 10\% discount on all purchases.}

\sense{vt.}{打折出售}
\ex{The shop discounted the damaged goods.}

\begin{origin}
来自古法语 \textit{desconter}（扣减）：\textit{dis-}（减）+ \textit{computare}（计算）——compute（计算）的同源词。同源词：\textbf{count}、\textbf{account}、\textbf{compute}。"算掉一部分"就是折扣。
\end{origin}

\end{document}
